\documentclass{article}
\usepackage[utf8]{inputenc}
\usepackage[T1]{fontenc}
\usepackage[margin=2cm]{geometry}
\usepackage{hyperref}
\sloppy
\begin{document}
\section*{Historical source}
Encyclopaedia Britannica scan, source row 1580. Parse 3.5 extracted the text. Super 3.5 VL corrected it from an enlarged crop of each block. It can still contain recognition errors.
Page headers, footers, and page numbers are removed.

\subsection*{TRIGONOMETRY (E0131)}
TRIGONOMETRY.

\textbackslash{}begin\{tabular\}\{cc\}
b to their difference 21 \& 1:32222\textbackslash{}\textbackslash{}
So is the tangent of half the fum of? \& 84 \& 1.92428\textbackslash{}\textbackslash{}
the angle ID and C \& 610 58' 10.27372\textbackslash{}\textbackslash{}
To the tangent of half their difference 79 50' \& 9.13806\textbackslash{}\textbackslash{}
So by theorem, we have the angles D and C thus; \& \textbackslash{}\textbackslash{}
To hold the fun of the angles I D and C \& 611 58'\textbackslash{}\textbackslash{}
Add half their difference \& 70 50'\textbackslash{}\textbackslash{}
And the fum is the greater angle D \& 69o 48'\textbackslash{}\textbackslash{}
Alfo, from half the fum \& 61o 58'\textbackslash{}\textbackslash{}
Take half the difference \& 70 50'\textbackslash{}\textbackslash{}
And there remains the lifter angle C \& 54°08'\textbackslash{}\textbackslash{}
Then by Cafel II, we have the following analogy for find ing On the leg required, viz. \& S.C. B D :: S, B. D C.\textbackslash{}\textbackslash{}
I. As the fine of C \& 54°03' 9.90860\textbackslash{}\textbackslash{}
To BD \& 133 \& 2.12385\textbackslash{}\textbackslash{}
So is the fine of B \& 56°03' 9.91883\textbackslash{}\textbackslash{}
TOO \& 150A 1.13590\textbackslash{}\textbackslash{}
Case VI. Three files being given, to find the angles. \& \textbackslash{}\textbackslash{}
Examle: In the triangle A B C ((bit a 9) ) fuppfe AB=13, AC=185 7, and BC=B4 y required the angles \& \textbackslash{}\textbackslash{}
A, B, and C; To the triangle will be confucted, and the angles may be incurred by a line of clouds. \& \textbackslash{}\textbackslash{}
II. By undertaking: Let fall the perpendicular, BD from the vertex B, upon the bafe’AC; which will divide the one into two segments AD and DC, the lengths where of may be found by theorem 4, thus: \& \textbackslash{}\textbackslash{}
As the bafe AC \& 185 7 \& 2.26893\textbackslash{}\textbackslash{}
To the form of the sides AB and BC 240 \& 2.8034\textbackslash{}\textbackslash{}
So is the difference of the files 72 \& 1.85733\textbackslash{}\textbackslash{}
To the diff, of the figures of the bafa 93 \& 1.0687i\textbackslash{}\textbackslash{}
And having the fum of the fragments, viz. the whole bafs, and their difference, we find the fragments themselves, by theorem 3, thus: \& \textbackslash{}\textbackslash{}
To half the fum of the fragments \& 92E-6\textbackslash{}\textbackslash{}
And half their difference \& 46-5\textbackslash{}\textbackslash{}
And the fum is the greater fragment AD \& 139-7\textbackslash{}\textbackslash{}
Allo from half the fum of the fragment \& 52-8\textbackslash{}\textbackslash{}
Take half their difference \& 46-5\textbackslash{}\textbackslash{}
The remainder is the lifter fragment DC \& 46-3\textbackslash{}\textbackslash{}
Now the triangle ABC is divided, by the perpendicular DB, into two right angles triangles, ADB and DRG; in the lift of which are given the hypothecule AB=156, and the bafe AD=130-3, to find the oblique angles, for which we have (by Cafel V, of rechargular trigonometry) the fol-lowing analogy, viz. \& \textbackslash{}\textbackslash{}
As AB \& 159 \& 2.10312\textbackslash{}\textbackslash{}
Tier AD \& 193 \& 2.1359\textbackslash{}\textbackslash{}
So is the radius \& 90° \& 1.00000\textbackslash{}\textbackslash{}
To the co-fise of the angle A \& 26° 40° \& 9.95083\textbackslash{}\textbackslash{}
\textbackslash{}end\{tabular\}

Alfo the angle C is found by the fam cafe, thus:\textbackslash{}\textbackslash{}
As IB 84 \& 1.92428\textbackslash{}\textbackslash{}
To CD \& 46:3 \& 1.66518\textbackslash{}\textbackslash{}
So is the radius \& 90° \& 10.0000\textbackslash{}\textbackslash{}
To the co-fine of C \& 56° 30° \& 9.74130\textbackslash{}\textbackslash{}
Having found the two angles A and C, we have the third, B, by taking the time of the other two from 150, thus: \& \textbackslash{}\textbackslash{}
The fum of all the three angles is 180° \& The fum of A and C is \& 83° 10°\textbackslash{}\textbackslash{}
The angle B is \& 96° 50° \& \textbackslash{}\textbackslash{}
All the proportions ufd for the foliations of the fever et cafes in plain trigonometry, may be performed by the final and compd. On the leaf there are several logarithmic lines, viz. one of numbers, another of fines, and one of tangents, éx. \& \textbackslash{}\textbackslash{}
And the way of working a proportion by these is this, viz., extend your compaifles from the frit term of your pro portion, found on the feale, to the second; and with that extent, flying on foot in the third term, the other will reach the fourth term required. \& \textbackslash{}\textbackslash{}
SPHERICAL TRIGONOMETRY \& \textbackslash{}\textbackslash{}
\textbackslash{}end\{tabular\}

SPHERICAL TRIGONOMETRY is the art whereby, from three given parts of a spherical triangle, we discover the rest; and, like plane trigonometry, is either right-angled, or oblique-angled. But before we give the analogies for the solution of the several cases in either, it will be proper to premise the following theorems.

Theorem I. In all right-angled spherical triangles, the fine of the hypotenuse: radius :: fine of a leg : fine of its opposite angle. And the fine of a leg : radius :: tangent of the other leg : tangent of its opposite angle.

\subsection*{DEMONSTRATION (E0132)}
DEMONSTRATION. Let E D A F G (ibid. fig. 3.) represent the eighth part of a Sphere, where the quadrantal planes EDFG, EDBC, are both perpendicular to the quadrantal plane ADFB; and the quadrantal plane ADGC is perpendicular to the plane EDFG; and the spherical triangle ABC is right-angled at B, where CA is the hypotenuse, and BA, BC, are the legs.

To the arches GF, CB, draw the tangents HF, OB, and the fines GM, CI on the radii DF, DB; also draw BL the fine of the arch AB, and CK the fine of AC; and then join IK and OL. Now HF, OB, GM, CI, are all perpendicular to the plane ADFB. And HD, GK, OL, lie all in the same plane ADGC. Also FD, IK, BL, lie all in the same plane ADGC. Therefore, the right-angled triangles HFD, CIK, ODL, having the equal angles HDF, CKI, OLB, are similar. And CK : DG :: CI : GM; that is, as the fine of the hypotenuse : rad. :: fine of a leg : fine of its opposite angle. For GM is the fine of the arc GF, which measures the angle CAB. Also, LB : DF :: BO : SH; that is, as the fine of a leg : radius :: tangent of the other leg : tangent of its opposite angle. Q. E. D. Hence it follows, that the fines of the angles of any oblique spherical triangle ACD (ibid. n° 2.) are to one another, directly, as the fines of the opposite sides. Hence it also follows, that, in right-angled spherical triangles, having the same perpendicular, the fines of the bases will be to each other, inversely, as the tangents of the angles at the bases.


\section*{Source record}
\begin{description}
\item[Scan ID] \texttt{britannica:309f765b1ecc83e0a7f7762e44b4a927a7555443:train:1580}
\item[Image file] \texttt{page-1580-e0bdcd3f241e.jpg}
\item[Image SHA-256] \texttt{e0bdcd3f241e27ff8e01ab1d9d2d473bc5373e9ac54cdbd1887cf0407fbf0eaa}
\item[Dataset] NationalLibraryOfScotland/encyclopaedia\_britannica\_illustrated at revision \texttt{309f765b1ecc83e0a7f7762e44b4a927a7555443} (CC0-1.0)
\item[Extraction] model \texttt{nvidia/nemotron-3.5-parse}, raw reply \texttt{extract/raw/page-1580-e0bdcd3f241e.json}
\item[Raw reply SHA-256] \texttt{8b443e8528941e6ee5a18e76fdb489e1c7b1c8ae9dfcadcec4a735beb1dbf97a}
\item[Correction] model \texttt{nvidia/Nemotron-3.5-Super-VL}, block crops, raw replies \texttt{correct/raw/crop\_nothink/page-1580-e0bdcd3f241e-b*.json}; parse: 2, super\_vl: 5
\item[Running header] none
\end{description}
\end{document}
